% =====================================================================
% 15. Introduction
%
% The text is the Introduction section of the paper, word for word.
% Only the outline of the thesis is added at the end, as the paper has
% no counterpart for it.
%
% The multipliers of the last paragraph are left as the red
% placeholders of item 13, because the numbers the paper reports here
% disagree with the ones of the Evaluation and of the Conclusions.
% TODO: fill them once the final numbers are in.
% =====================================================================
\chapter{Introduction}
\label{ch:introduction}

Neural network classifiers are highly accurate but vulnerable to
adversarial examples: small input perturbations that change their
classification~\cite{szegedy2014,carlini2017,madry2018}. To provide
guarantees against them, \emph{local robustness} verification
certifies that a single given input cannot be misclassified under any
small
perturbation~\cite{katz2017reluplex,gehr2018,singh2019,wang2021}. But
a deployed network in a safety-critical application, such as
autonomous driving~\cite{eykholt2018} or medical
imaging~\cite{finlayson2019}, faces inputs it has never seen, and must
know which of them an attacker could flip. \emph{Global robustness}
addresses this: it guarantees that any input classified above a
confidence level cannot be flipped by a
perturbation~\cite{leino2021gloro,kabaha2024vhagar,athavale2024}. At
inference, any input below this level is flagged as potentially
vulnerable and can be handled accordingly.

Verifying global robustness, however, is hard: it reasons about a
perturbation of every input, which requires analyzing two copies of
the network, the clean input and its perturbation, rather than the
single network of local robustness, which doubles the complexity.
Existing verifiers thus face a tradeoff: exact
verification~\cite{kabaha2024vhagar} takes hours even on networks with
up to 2000 neurons, while over-approximating
it~\cite{athavale2024} loses precision and proves robustness for fewer
inputs.

We observe that lightweight abstractions such as intervals, zonotopes,
and triangle relaxations made local robustness
tractable~\cite{liu2021algorithms}. Extending them to global
robustness naively loses too much precision, and using each on its own
yields only a moderate speedup. Our key idea is to exploit the
relationship between the two copies with abstractions that together
keep the over-approximation error small while speeding up the
verification.

We introduce \tool (Boosting by Lightweight Encoding of Network
Differences), which extends these abstractions to the two network
copies: a perturbation-difference interval bounds how much the two
copies can differ at each neuron, extending difference intervals for
various perturbations (including linear transformations), a zonotope
tightens the common per-neuron interval bounds by capturing
dependencies, and a novel conditional triangle linearly relaxes
neurons whose activation is nearly determined, leveraging the relation
between the two copies. \tool is also the first to support incremental
global robustness verification: if a verified network is
retrained~\cite{parisi2019} or fine-tuned~\cite{yosinski2014}, \tool
verifies the revised network by transferring the old bounds and
warm-starting the analysis.

We evaluate \tool on MNIST~\cite{lecun1998}, Fashion-MNIST~\cite{xiao2017},
CIFAR-10~\cite{krizhevsky2009} and HAR~\cite{anguita2013har} networks
across seven perturbation types. \tool speeds up the exact verifier
\baseline~\cite{kabaha2024vhagar} by $7.6$ times on
average (up to $102.6$ times), and on the hardest tasks,
where both time out, returns bounds $19.9$ times tighter
(up to $261.2$ times). For incremental verification, \tool
speeds up verification of a revised network by $12.5$
times (up to $115.4$ times) and returns
$20.9$ times tighter bounds (up to $261.2$
times) on the tasks that time out.


\section*{Thesis Organization}

Chapter~\ref{ch:background} defines neural network classifiers,
global robustness and its exact verification.
Chapter~\ref{ch:related-work} reviews the methods that verify
robustness and the limitations that leave global robustness open.
Chapter~\ref{ch:key-ideas} presents our key ideas, 
and Chapter~\ref{ch:blend} defines
them formally.
Chapter~\ref{ch:evaluation} evaluates \tool against the previous 
state-of-the-art verifier, in execution time and in the tightness of the
bounds, and measures the contribution of each of our techniques.
Chapter~\ref{ch:discussion} concludes and presents the directions we
leave open.
